\documentclass[11pt, a4paper]{report}

% ========== Common Packages ==========
\usepackage{fontspec}
\usepackage{kotex}
\usepackage{amsmath, amssymb, amsthm}
\usepackage{mathtools}
\usepackage{bm}
\usepackage{graphicx}
\usepackage{booktabs}
\usepackage{array}
\usepackage{tabularx}
\usepackage{longtable}
\usepackage{hyperref}
\usepackage{cleveref}
\usepackage{geometry}
\usepackage{fancyhdr}
\usepackage{titlesec}
\usepackage{enumitem}
\usepackage{xcolor}
\usepackage{tcolorbox}
\usepackage{algorithm}
\usepackage{algorithmic}
\usepackage{tikz}
\usetikzlibrary{arrows.meta, positioning, calc, decorations.pathreplacing, shapes.geometric, fit, patterns, angles, quotes, trees}
\usepackage{pgfplots}
\pgfplotsset{compat=1.18}
\usepackage{caption}
\usepackage{subcaption}
\usepackage{float}

\geometry{margin=2.5cm, top=3cm, bottom=3cm}

\usepackage{multirow}
% ========== Colors ==========
% Common
\definecolor{defblue}{RGB}{30, 80, 150}
\definecolor{thmgreen}{RGB}{20, 110, 60}
\definecolor{noteorange}{RGB}{200, 100, 0}
\definecolor{lightblue}{RGB}{230, 240, 255}
\definecolor{lightgreen}{RGB}{230, 255, 235}
\definecolor{lightorange}{RGB}{255, 245, 230}
\definecolor{lightgray}{RGB}{245, 245, 245}
% Extended
\definecolor{lightred}{RGB}{255, 235, 235}
\definecolor{darkred}{RGB}{180, 30, 30}
\definecolor{biogreen}{RGB}{10, 130, 80}
\definecolor{cvpurple}{RGB}{100, 40, 150}
\definecolor{injuryred}{RGB}{200, 40, 40}
\definecolor{codegray}{RGB}{240, 240, 245}
\definecolor{pipeblue}{RGB}{25, 100, 180}
\definecolor{constraintred}{RGB}{200, 50, 50}
\definecolor{termblack}{RGB}{30, 30, 30}
\definecolor{goldcolor}{RGB}{180, 140, 20}


% ========== Korean theorem names (이 파일 전용) ==========
\theoremstyle{definition}
\newtheorem{definition}{정의}[chapter]
\newtheorem{theorem}{정리}[chapter]
\newtheorem{proposition}{명제}[chapter]
\theoremstyle{remark}
\newtheorem{remark}{참고}[chapter]

% ========== Korean tcolorbox (이 파일 전용) ==========
\newtcolorbox{keyidea}[1][]{colback=lightblue, colframe=defblue, fonttitle=\bfseries, title={핵심}, #1}
\newtcolorbox{importantnote}[1][]{colback=lightorange, colframe=noteorange, fonttitle=\bfseries, title={중요}, #1}
\newtcolorbox{warningbox}[1][]{colback=lightred, colframe=darkred, fonttitle=\bfseries, title={주의}, #1}
\newtcolorbox{proofbox}[1][]{colback=lightgreen, colframe=thmgreen, fonttitle=\bfseries, title={유도}, #1}
\newtcolorbox{practicebox}[1][]{colback=lightgray, colframe=gray!60, fonttitle=\bfseries, title={우리 프로젝트}, #1}

% ========== Custom math commands ==========
% Vectors (lowercase)
\newcommand{\vx}{\mathbf{x}}
\newcommand{\vd}{\mathbf{d}}
\newcommand{\vo}{\mathbf{o}}
\newcommand{\vc}{\mathbf{c}}
\newcommand{\vr}{\mathbf{r}}
\newcommand{\vp}{\mathbf{p}}
\newcommand{\vv}{\mathbf{v}}
\newcommand{\vn}{\mathbf{n}}
\newcommand{\vu}{\mathbf{u}}
\newcommand{\vt}{\mathbf{t}}
\newcommand{\ve}{\mathbf{e}}
\newcommand{\vq}{\mathbf{q}}
% Vectors (uppercase)
\newcommand{\vF}{\mathbf{F}}
\newcommand{\vM}{\mathbf{M}}
\newcommand{\va}{\mathbf{a}}
\newcommand{\vg}{\mathbf{g}}
\newcommand{\vX}{\mathbf{X}}
% Bold Greek
\newcommand{\vmu}{\bm{\mu}}
\newcommand{\vbeta}{\bm{\beta}}
\newcommand{\vtheta}{\bm{\theta}}
\newcommand{\vpsi}{\bm{\psi}}
% Matrices
\newcommand{\mSigma}{\bm{\Sigma}}
\newcommand{\mR}{\mathbf{R}}
\newcommand{\mS}{\mathbf{S}}
\newcommand{\mW}{\mathbf{W}}
\newcommand{\mJ}{\mathbf{J}}
\newcommand{\mG}{\mathbf{G}}
\newcommand{\mT}{\mathbf{T}}
\newcommand{\mI}{\mathbf{I}}
\newcommand{\mB}{\mathbf{B}}
\newcommand{\mK}{\mathbf{K}}
\newcommand{\mP}{\mathbf{P}}
\newcommand{\mF}{\mathbf{F}}
\newcommand{\mE}{\mathbf{E}}
\newcommand{\mA}{\mathbf{A}}
\newcommand{\mH}{\mathbf{H}}
% Sets and groups
\newcommand{\RR}{\mathbb{R}}
\newcommand{\SO}{\mathrm{SO}}
% Units
\newcommand{\degps}{\,\text{deg/s}}
\newcommand{\Nm}{\ensuremath{\,\text{N}\cdot\text{m}}}
\newcommand{\BW}{\,\text{BW}}

% ========== Header/Footer ==========
\pagestyle{fancy}
\fancyhf{}
\fancyhead[L]{\leftmark}
\fancyhead[R]{\thepage}
\renewcommand{\headrulewidth}{0.4pt}

% ========== Title formatting ==========
\titleformat{\chapter}[display]
  {\normalfont\huge\bfseries\color{defblue}}
  {\chaptertitlename\ \thechapter}{20pt}{\Huge}

\titleformat{\section}
  {\normalfont\Large\bfseries\color{defblue!80!black}}
  {\thesection}{1em}{}

\titleformat{\subsection}
  {\normalfont\large\bfseries\color{defblue!60!black}}
  {\thesubsection}{1em}{}


\begin{document}

\begin{titlepage}
\centering
\vspace*{1.5cm}
{\Large\color{gray} 수학의 정석 시리즈}
\vspace{0.5cm}

{\Huge\bfseries\color{defblue} COLMAP은 무엇이고\\[0.3cm]
왜 필요한가}

\vspace{0.8cm}

{\LARGE\color{defblue!70!black} 카메라로 사진을 찍으면 왜 COLMAP이 나오고,\\
그것이 3DGS와 SMPL 피팅에\\
어떻게 연결되는가}

\vspace{3cm}
{\large 2026}

\vspace{\fill}
{\small\color{gray} ``카메라 보정이 왜 필요한가''부터\\
``COLMAP 출력이 SMPL 피팅의 어디에 쓰이는가''까지 완전 해설.}
\end{titlepage}

\tableofcontents
\newpage


% ====================================================================
\chapter{카메라로 사진을 찍으면 무슨 일이 일어나는가}
% ====================================================================

\section{사진 = 3D 세계의 2D 투영}

카메라로 사진을 찍는 것은 수학적으로 \textbf{3D $\rightarrow$ 2D 투영}이다.

\begin{figure}[H]
\centering
\begin{tikzpicture}[scale=0.9]
% 3D world
\draw[thick, defblue] (6, 0) -- (6, 3) -- (6.5, 3.5) -- (6, 3) -- (5.5, 3.5);
\draw[thick, defblue] (6, 0) -- (5.5, -0.5);
\draw[thick, defblue] (6, 0) -- (6.5, -0.5);
\fill[defblue] (6, 3.7) circle (0.2);
\node[font=\small, text=defblue] at (6, -1.2) {3D 세계의 사람};

% Camera
\fill[gray!30] (0, 1.5) -- (-0.3, 1.8) -- (0.3, 1.8) -- cycle;
\fill[gray!30] (-0.3, 1.2) rectangle (0.3, 1.8);
\node[font=\small] at (0, 0.8) {카메라};

% Image plane
\draw[thick, noteorange] (2, 0) -- (2, 3);
\node[font=\small, text=noteorange] at (2, -0.5) {이미지 (2D)};

% Projection lines
\draw[gray, dashed] (0, 1.5) -- (6, 3.7);
\draw[gray, dashed] (0, 1.5) -- (6, 0);
\draw[gray, dashed] (0, 1.5) -- (6, -0.5);

% Projected person on image
\draw[thick, noteorange] (2, 1.0) -- (2, 2.2);
\fill[noteorange] (2, 2.35) circle (0.08);

% Labels
\draw[<->, thick] (0.5, 1.5) -- node[above, font=\tiny] {$f$ (초점거리)} (1.8, 1.5);
\draw[<->, thick] (2.2, 0) -- node[right, font=\tiny] {이미지 높이 $H$} (2.2, 3);
\end{tikzpicture}
\caption{카메라는 3D 세계를 2D 이미지로 투영한다. 이 투영 과정을 수학적으로 기술한 것이 \textbf{카메라 모델}이다.}
\end{figure}

사진 한 장에서 알 수 있는 것: 사람의 \textbf{2D 위치, 크기, 모양}.

사진 한 장에서 알 수 \textbf{없는} 것: 사람이 카메라에서 \textbf{얼마나 먼지} (깊이).

\begin{keyidea}
하나의 2D 이미지에서는 깊이 정보가 손실된다.
\textbf{여러 대의 카메라}에서 같은 장면을 찍으면, 각 카메라의 시차(parallax)로 깊이를 복원할 수 있다.
이것이 \textbf{다시점 3D 재구성}의 원리이며, COLMAP이 하는 일이다.
\end{keyidea}


\section{투영을 되돌리려면: 두 가지를 알아야 한다}

3D $\rightarrow$ 2D 투영을 거꾸로 (2D $\rightarrow$ 3D) 하려면 두 가지 정보가 필요하다:

\begin{enumerate}[leftmargin=*]
  \item \textbf{카메라 내부 파라미터 $\mK$} (intrinsics): 렌즈의 초점 거리, 이미지 중심 위치.
  ``이 카메라가 3D를 2D로 \textbf{어떻게} 변환하는가''

  \item \textbf{카메라 외부 파라미터 $[\mR|\vt]$} (extrinsics): 카메라의 위치와 방향.
  ``이 카메라가 3D 세계에서 \textbf{어디에 있고 어디를 보는가}''
\end{enumerate}

\begin{definition}[카메라 보정 (Camera Calibration)]
카메라 내부 파라미터 $\mK$와 외부 파라미터 $[\mR|\vt]$를 추정하는 과정.
이것이 \textbf{COLMAP이 하는 핵심 작업}이다.
\end{definition}


% ====================================================================
\chapter{COLMAP은 무엇인가}
% ====================================================================

\section{정의}

\begin{definition}[COLMAP]
COLMAP은 여러 장의 2D 사진으로부터:
\begin{enumerate}[nosep]
  \item 각 카메라의 내부 파라미터 $\mK$ (초점 거리, 주점)
  \item 각 카메라의 외부 파라미터 $[\mR|\vt]$ (위치, 방향)
  \item 장면의 3D 점군 (sparse point cloud)
\end{enumerate}
을 자동으로 추정하는 소프트웨어이다.
이 과정을 \textbf{Structure from Motion (SfM)}이라 한다.
\end{definition}


\section{COLMAP이 작동하는 원리}

\begin{figure}[H]
\centering
\begin{tikzpicture}[
  block/.style={rectangle, rounded corners=4pt, draw=defblue, fill=lightblue,
    minimum width=3cm, minimum height=0.8cm, font=\small\bfseries, align=center},
  arrow/.style={-Stealth, thick, defblue!60}
]
\node[block] (feat) at (0,0) {1. 특징점 검출\\(SIFT)};
\node[block] (match) at (5,0) {2. 특징점 매칭};
\node[block] (init) at (0,-2) {3. 초기 2장\\카메라 추정};
\node[block] (add) at (5,-2) {4. 카메라 추가\\(증분적)};
\node[block] (ba) at (2.5,-4) {5. Bundle\\Adjustment};
\node[block] (out) at (2.5,-6) {$\mK, [\mR|\vt]$\\+ 3D 점군};

\draw[arrow] (feat) -- (match);
\draw[arrow] (match) -- (init);
\draw[arrow] (init) -- (add);
\draw[arrow] (add) -- (ba);
\draw[arrow, dashed] (ba) -- (add) node[midway, right, font=\tiny] {반복};
\draw[arrow] (ba) -- (out);
\end{tikzpicture}
\caption{COLMAP SfM 파이프라인}
\end{figure}

\subsection{Step 1: 특징점 검출 (Feature Extraction)}

각 이미지에서 \textbf{독특한 점}들을 찾는다. COLMAP은 SIFT (Scale-Invariant Feature Transform)를 사용한다.

특징점이란: 코너, 가장자리, 텍스처가 풍부한 점 등 다른 이미지에서도 같은 점을 식별할 수 있는 위치.

\begin{practicebox}
우리 데이터에서 이미지당 $\sim$5,800개의 SIFT 특징점이 검출되었다 (충분한 양).
\end{practicebox}


\subsection{Step 2: 특징점 매칭 (Feature Matching)}

서로 다른 이미지에서 \textbf{같은 3D 점}을 보고 있는 특징점 쌍을 찾는다.

이미지 A의 특징점 $\mathbf{f}_i^A$와 이미지 B의 특징점 $\mathbf{f}_j^B$가 같은 3D 점이면,
이 둘은 ``매치''이다.

\begin{importantnote}
매칭이 정확하려면, \textbf{같은 3D 점이 여러 이미지에서 같은 위치에 있어야 한다}.
이것이 COLMAP이 \textbf{정적 장면}을 가정하는 이유이다.
움직이는 물체의 특징점은 프레임마다 다른 3D 위치에 있으므로 매칭이 오염된다.
\end{importantnote}


\subsection{Step 3--4: 카메라 추정 (증분적 SfM)}

매칭된 특징점 쌍으로부터 카메라 파라미터를 추정한다.

\textbf{2장에서 시작}: 가장 많은 매칭을 가진 이미지 쌍으로 초기 카메라 2대를 추정.

\textbf{1장씩 추가}: 이미 추정된 3D 점과의 대응을 통해 새 카메라의 위치를 계산 (PnP 문제).

이 과정이 $N$장의 이미지 모두가 등록될 때까지 반복된다.


\subsection{Step 5: Bundle Adjustment}

모든 카메라 파라미터와 3D 점 위치를 \textbf{동시에 정밀화}하는 비선형 최적화:

\begin{equation}
\min_{\{\mK_c, \mR_c, \vt_c\}, \{\vX_i\}}
\sum_{c} \sum_{i \in \mathcal{V}_c} \left\| \vu_{ci} - \pi_c(\vX_i) \right\|^2
\label{eq:ba}
\end{equation}

여기서:
\begin{itemize}[nosep]
  \item $\vu_{ci}$: 카메라 $c$에서 3D 점 $i$의 검출된 2D 위치
  \item $\pi_c(\vX_i) = \mK_c [\mR_c | \vt_c] \tilde{\vX}_i$: 추정된 파라미터로 계산한 2D 투영
  \item $\mathcal{V}_c$: 카메라 $c$에서 보이는 3D 점의 집합
\end{itemize}

\begin{keyidea}
Bundle Adjustment는 ``모든 카메라에서 모든 3D 점의 재투영 오차를 최소화''한다.
이것이 COLMAP 출력의 정확도를 결정하는 핵심 단계이다.
\end{keyidea}


\section{COLMAP의 출력}

COLMAP이 최종적으로 출력하는 것:

\begin{table}[H]
\centering
\caption{COLMAP 출력 파일}
\renewcommand{\arraystretch}{1.3}
\begin{tabularx}{\textwidth}{llX}
\toprule
\textbf{파일} & \textbf{내용} & \textbf{수학적 의미} \\
\midrule
\texttt{cameras.txt} & 카메라 모델 + 내부 파라미터 & $\mK = \begin{pmatrix} f_x & 0 & c_x \\ 0 & f_y & c_y \\ 0 & 0 & 1 \end{pmatrix}$ \\
\texttt{images.txt} & 이미지별 외부 파라미터 & $[\mR | \vt]$ (쿼터니언 + 이동 벡터) \\
\texttt{points3D.txt} & 3D 점군 & $\{\vX_i\} \subset \RR^3$ + RGB 색상 \\
\bottomrule
\end{tabularx}
\end{table}

\begin{keyidea}
COLMAP 출력의 $\mK$는 \textbf{입력 이미지의 픽셀 좌표}와 직접 대응한다.
1920$\times$1080 이미지를 넣으면, $c_x \approx 960$, $c_y \approx 540$으로 나온다.
\textbf{좌표 변환이 필요 없다} --- 이것이 COLMAP의 가장 큰 장점이다.
\end{keyidea}


% ====================================================================
\chapter{COLMAP이 우리 프로젝트에서 왜 필요한가}
% ====================================================================

\section{전체 파이프라인에서의 위치}

\begin{figure}[H]
\centering
\begin{tikzpicture}[
  block/.style={rectangle, rounded corners=3pt, draw, minimum width=2.8cm, minimum height=0.7cm,
    align=center, font=\small},
  arrow/.style={-{Stealth[length=2mm]}, thick}
]
\node[block, fill=lightblue, draw=defblue, very thick] (colmap) at (0,0) {\textbf{COLMAP}\\$\mK, [\mR|\vt]$};
\node[block, fill=lightgray] (op2d) at (4,0) {OpenPose 2D\\$\vu_{ck}$};
\node[block, fill=cvpurple!10, draw=cvpurple] (smpl) at (2,-2) {EasyMocap\\SMPL 피팅};
\node[block, fill=lightgreen, draw=thmgreen] (gart) at (2,-4) {GART\\3DGS 학습};

\draw[arrow, defblue] (colmap) -- (smpl) node[midway, left, font=\tiny] {$\pi_c(\cdot)$에 사용};
\draw[arrow] (op2d) -- (smpl) node[midway, right, font=\tiny] {$\vu_{ck}$};
\draw[arrow] (smpl) -- (gart) node[midway, right, font=\tiny] {$\vtheta, \vbeta$};
\draw[arrow, defblue] (colmap) -- (gart) node[midway, left, font=\tiny] {렌더링에 사용};
\end{tikzpicture}
\caption{COLMAP 출력($\mK, [\mR|\vt]$)은 SMPL 피팅과 GART 학습 \textbf{모두}에 필요하다.}
\end{figure}

\subsection{SMPL 피팅에서의 역할}

EasyMocap의 목적 함수:
\begin{equation}
\min_{\vtheta, \vbeta} \sum_{c=1}^{7} \sum_{k} w_{ck} \left\| \vu_{ck} - \underbrace{\mK_c [\mR_c | \vt_c]}_{\text{COLMAP 출력}} J_k(\vtheta, \vbeta) \right\|^2
\end{equation}

$\mK_c$와 $[\mR_c|\vt_c]$가 없으면 $\pi_c(\cdot)$를 계산할 수 없고,
SMPL의 3D 관절을 2D로 투영할 수 없다.
\textbf{SMPL 피팅 자체가 불가능하다.}


\subsection{3DGS/GART 학습에서의 역할}

3DGS는 각 카메라에서 본 이미지를 렌더링하여 실제 이미지와 비교한다:
\begin{equation}
\mathcal{L} = \sum_c \left\| I_c^{\text{rendered}}(\mK_c, \mR_c, \vt_c) - I_c^{\text{real}} \right\|
\end{equation}

렌더링에 $\mK_c$와 $[\mR_c|\vt_c]$가 필요하다.
\textbf{카메라 파라미터 없이는 3DGS도 불가능하다.}


\section{왜 COLMAP이 실패했는가}

\begin{warningbox}
COLMAP은 \textbf{정적 장면}을 가정한다 (Bundle Adjustment 수식 \cref{eq:ba} 참조).

우리 데이터에서 실패한 이유:
\begin{enumerate}[nosep]
  \item \textbf{투수가 움직임}: 사람 위의 SIFT 특징점이 프레임마다 다른 3D 위치.
  COLMAP은 이 점들이 ``정적''이라 가정하고 3D 위치를 추정하므로, 결과가 오염됨.

  \item \textbf{7대 카메라가 동일 기종}: COLMAP이 인식 못해 서로 다른 $f_x$를 추정.
  GT: 모든 카메라 $f \approx 790$\,px. COLMAP: 531--1367\,px (최대 2.6배 차이).

  \item \textbf{일부 카메라 시야 겹침 부족}: cam4 등록률 8\% (Round 1).
\end{enumerate}
\end{warningbox}


\section{해결 방법: 마스킹 후 COLMAP}

\begin{keyidea}
\textbf{사람을 이미지에서 제거}하면 COLMAP의 정적 장면 가정이 만족된다.

SAM3 마스크가 이미 있으므로:
\begin{enumerate}[nosep]
  \item 원본 이미지에서 사람 영역(마스크)을 검은색으로 채움
  \item 배경만 남은 이미지로 COLMAP 실행
  \item 출력 $\mK$는 \textbf{원본 해상도 기준} --- 좌표 변환 불필요
  \item OpenPose 키포인트($\vu_{ck}$)와 바로 호환
\end{enumerate}

이것이 DA3 역변환 문제를 \textbf{완전히 우회}하는 방법이다.
\end{keyidea}


% ====================================================================
\chapter{COLMAP vs DA3: 왜 DA3에서 어려웠는가}
% ====================================================================

\section{근본 차이}

\begin{table}[H]
\centering
\caption{COLMAP vs DA3 카메라 보정 비교}
\renewcommand{\arraystretch}{1.4}
\begin{tabularx}{\textwidth}{lXX}
\toprule
& \textbf{COLMAP} & \textbf{DA3} \\
\midrule
원리 & 특징점 매칭 + 기하 & 딥러닝 (깊이 + 포즈) \\
입력 해상도 & \textbf{원본 그대로} & 내부 리사이즈/크롭 \\
출력 $\mK$ 기준 & \textbf{입력 이미지 해상도} & DA3 처리 해상도 \\
좌표 변환 & \textbf{불필요} & 역변환 필수 \\
정적 장면 가정 & \textbf{필요} (움직이면 실패) & 불필요 \\
동적 장면 & $\times$ & $\bigcirc$ (이론적) \\
\bottomrule
\end{tabularx}
\end{table}

\section{DA3에서 발생한 좌표 변환 문제}

\begin{figure}[H]
\centering
\begin{tikzpicture}[
  img/.style={rectangle, draw, minimum width=2cm, minimum height=1.2cm, align=center, font=\tiny},
  arrow/.style={-Stealth, thick},
  label/.style={font=\scriptsize\color{darkred}}
]
% Original
\node[img, minimum width=3cm, minimum height=1.8cm, fill=lightblue!30] (orig) at (0,0) {1920$\times$1080\\원본};

% DA3 resize
\node[img, minimum width=1.5cm, minimum height=0.9cm, fill=lightorange!30] (resized) at (4,0) {504$\times$280\\리사이즈};

% DA3 crop
\node[img, minimum width=0.9cm, minimum height=0.9cm, fill=lightred!30] (cropped) at (7,0) {280$\times$280\\크롭};

% K labels
\node[label] at (0, -1.5) {$\mK_{\text{orig}}$: $c_x \approx 960$};
\node[label] at (4, -1.5) {$\mK_{\text{resized}}$: $c_x \approx 252$};
\node[label] at (7, -1.5) {$\mK_{\text{cropped}}$: $c_x \approx 140$};

\draw[arrow] (orig) -- (resized) node[midway, above, font=\tiny] {$s_x, s_y$};
\draw[arrow] (resized) -- (cropped) node[midway, above, font=\tiny] {$-112$};

% OpenPose
\node[img, minimum width=3cm, minimum height=1.8cm, fill=lightgreen!30] (op) at (0, -3.5) {OpenPose\\$\vu_{ck}$: 원본 좌표};

% Mismatch
\draw[darkred, very thick, dashed] (op) -- (cropped)
  node[midway, font=\small\bfseries\color{darkred}, fill=white] {좌표계 불일치!};

\end{tikzpicture}
\caption{DA3의 리사이즈+크롭으로 $\mK$의 좌표계가 변하지만, OpenPose 키포인트는 원본 좌표. 이 불일치가 모든 문제의 원인이었다.}
\end{figure}

\section{COLMAP이면 왜 해결되는가}

\begin{proofbox}[title={COLMAP의 좌표계 일관성}]
COLMAP에 1920$\times$1080 이미지를 입력하면:

\begin{enumerate}[nosep]
  \item COLMAP이 SIFT 특징점을 원본 해상도에서 검출: $(x_i, y_i)$ where $0 \leq x_i \leq 1920$
  \item Bundle Adjustment에서 이 좌표로 $\mK$를 최적화
  \item 출력 $\mK$의 $c_x \approx 960$, $c_y \approx 540$ (원본 해상도 기준)
  \item OpenPose도 같은 원본 이미지에서 키포인트 검출: $0 \leq u \leq 1920$
\end{enumerate}

$\Rightarrow$ COLMAP의 $\mK$와 OpenPose의 $\vu_{ck}$가 \textbf{자동으로 같은 좌표계}.

역변환, 스케일링, 패딩 오프셋 --- \textbf{모두 불필요}. $\square$
\end{proofbox}


% ====================================================================
\chapter{마스킹 후 COLMAP: 최종 전략}
% ====================================================================

\section{파이프라인}

\begin{figure}[H]
\centering
\begin{tikzpicture}[
  block/.style={rectangle, rounded corners=3pt, draw=defblue, fill=lightblue,
    minimum width=3.5cm, minimum height=0.7cm, align=center, font=\small},
  arrow/.style={-{Stealth[length=2mm]}, thick, defblue!60}
]
\node[block] (orig) at (0,0) {원본 이미지\\1920$\times$1080};
\node[block, fill=lightred!20, draw=darkred] (mask) at (5,0) {SAM3 마스크\\(이미 있음)};
\node[block] (masked) at (2.5,-1.5) {마스킹된 이미지\\(사람=검은색, 배경 유지)};
\node[block, very thick] (colmap) at (2.5,-3) {\textbf{COLMAP SfM}\\배경 특징점만 사용};
\node[block, fill=lightgreen, draw=thmgreen] (K) at (2.5,-4.5) {$\mK$ (원본 해상도)\\$[\mR|\vt]$};
\node[block] (op) at (7,-4.5) {OpenPose 2D\\$\vu_{ck}$ (원본 좌표)};
\node[block, fill=cvpurple!10, draw=cvpurple] (smpl) at (5,-6) {EasyMocap\\SMPL 피팅};

\draw[arrow] (orig) -- (masked);
\draw[arrow] (mask) -- (masked);
\draw[arrow] (masked) -- (colmap);
\draw[arrow] (colmap) -- (K);
\draw[arrow, thmgreen] (K) -- (smpl) node[midway, left, font=\tiny\color{thmgreen}] {좌표계 일치!};
\draw[arrow] (op) -- (smpl);
\end{tikzpicture}
\caption{마스킹 후 COLMAP $\rightarrow$ EasyMocap의 자연스러운 연결. 좌표 변환이 전혀 필요 없다.}
\end{figure}

\section{왜 이것이 올바른 해결책인가}

\begin{enumerate}[leftmargin=*]
  \item \textbf{정적 장면 가정 충족}: 배경은 움직이지 않으므로 COLMAP이 정상 작동
  \item \textbf{좌표계 자동 일치}: 입력 = 출력 = OpenPose = 모두 원본 해상도
  \item \textbf{역변환 불필요}: DA3의 리사이즈/크롭/패딩 문제 완전 회피
  \item \textbf{데이터 준비 완료}: SAM3 마스크 7대 $\times$ 1000프레임 이미 있음
  \item \textbf{카메라 공유 가능}: \texttt{single\_camera\_per\_folder}로 7대 카메라 내부 파라미터 공유
\end{enumerate}

\begin{importantnote}
이전 COLMAP 시도에서 \texttt{single\_camera=0} (이미지마다 독립 $\mK$)으로 설정하여
103개의 서로 다른 $\mK$가 추정되었다.

이번에는 \texttt{single\_camera\_per\_folder=1}로 설정하여
\textbf{7대 카메라에 대해 각각 1개의 $\mK$만} 추정하도록 한다.
또한 마스킹으로 사람이 제거되어 동적 장면 문제도 해결된다.
\end{importantnote}


\section{예상 결과}

마스킹 후 COLMAP이 성공하면:

\begin{table}[H]
\centering
\renewcommand{\arraystretch}{1.3}
\begin{tabular}{lll}
\toprule
\textbf{출력} & \textbf{예상 값} & \textbf{용도} \\
\midrule
$f_x \approx f_y$ & $\sim$790\,px (7대 일관) & SMPL 피팅 + GART \\
$c_x$ & $\sim$960\,px & 원본 해상도 기준 \\
$c_y$ & $\sim$540\,px & 원본 해상도 기준 \\
$[\mR|\vt]$ & 7대 카메라 상대 위치 & 삼각측량 + 렌더링 \\
재투영 오차 & $< 1$\,px & 정확도 검증 \\
\bottomrule
\end{tabular}
\end{table}


\end{document}
